\chapter{绪论}

\section{研究背景与意义}

随着互联网技术、物联网以及工业互联网的快速发展，网络协议已广泛应用于信息系统、嵌入式设备和关键基础设施之中。
协议作为系统之间进行信息交互和协同工作的基础，其实现的正确性和健壮性直接关系到系统的安全性与可靠性。
在实际应用中，协议实现往往涉及复杂的数据结构、严格的字段约束以及多阶段的状态转换逻辑，一旦在设计或实现过程中存在疏漏，极易引入安全缺陷。
攻击者可以通过构造异常报文触发程序中的边界条件错误或逻辑异常，从而造成服务拒绝、系统崩溃甚至更为严重的安全后果 \cite{de2015protocol, manes2019art}。
一个典型的例子是 2014 年爆发的 Heartbleed 漏洞 \cite{durumeric2014matter}，该漏洞源于 OpenSSL 实现 TLS/DTLS 协议心跳机制时边界检查缺失，导致攻击者能够远程泄露服务器内存中的敏感信息，造成了广泛的安全影响。


为了保障协议实现的安全性，研究人员和安全从业者一直在探索有效的测试方法来发现潜在的漏洞。
模糊测试（Fuzzing） \cite{manes2019art, ren2021survey} 是一种自动化的软件测试技术，其核心在于向计算机程序提供无效、非预期或随机的数据作为输入，以触发异常行为并发现潜在缺陷。
由于其无需完整源码支持、对测试目标适应性强，模糊测试在协议安全测试领域得到了广泛应用 \cite{zhang2024survey, li2011automated}。
在现有的模糊测试工具中，Peach\cite{gitlab_protocol_fuzzer_ce} 是一个具有代表性的经典框架，
其通过数据模型、状态模型、变异器和变异策略的组合，为复杂协议的自动化测试提供了良好的工程化基础，至今仍在工业界广泛应用。

Peach 框架提供了一套结构化的协议建模语言，允许用户定义协议字段的类型、长度、取值范围以及状态机等信息，以指导测试用例的生成和变异，其性能也在很大程度上受协议模型的构建质量影响\cite{zhang2023survey}。
然而，从协议规范文档中提取有效信息以构建高质量的协议模型是一项复杂而耗时的任务，
难以自动化完成，依赖测试人员手工编写，
尤其对于状态复杂、约束严格的协议而言更是如此。
而且，尽管可以定义出详细的数据模型和状态模型，在实际测试过程中，Peach 框架中的变异算子仍主要是字段类型相关的随机变异器，缺乏对特定协议语义的理解和利用。
此外，其结构化建模对于协议的字段间约束表达能力有限，只能描述例如长度约束、数量约束等基本关系，对字段之间更为复杂的语义关联和逻辑关系难以表达和利用。
这虽然能够基本保证测试用例在格式上合法，但随机变异往往破坏协议内部的语义一致性，导致大量生成的测试用例在解析或状态校验阶段即被丢弃，难以触发深层逻辑路径，制约了漏洞发现能力。

近年来，大语言模型（LLM）的兴起与飞速发展为协议模糊测试领域带来了新的机遇。
将大语言模型引入协议模糊测试过程 \cite{meng2024large, wang2025llmboofuzz}，
利用其在自然语言理解和知识抽取方面的显著优势，从协议规范文档中学习隐含规则和语义关系，
参与变异逻辑的生成，有望在保持输入结构合理性的同时，
更有针对性地对字段语义进行探索。

在此背景下，研究一种基于 Peach 模糊测试框架的协议语义感知变异增强方法，具有重要的理论意义和实践价值。该研究有助于弥补传统模糊测试框架在语义层面的不足，推动模糊测试技术向更加智能和高效的方向发展，同时也为复杂网络协议实现的安全测试提供新的思路。

\section{国内外研究现状}

\subsection{协议模糊测试}

网络协议模糊测试作为一种自动化提取协议实现漏洞的技术，其核心挑战在于处理协议数据的高度结构化语义以及复杂的有状态交互过程。如何在模糊测试的过程中有效融合协议语义感知与状态空间探索，构成了该领域研究的主要内容。

Peach\cite{gitlab_protocol_fuzzer_ce} 是一种典型的基于规约（Specification-based）的生成式模糊测试框架。它通过 XML 格式的“Peach Pit”文件定义协议的数据模型和状态模型。Peach 的优势在于能够根据预定义的字段类型、长度约束和校验和逻辑，生成符合协议语义规范的数据包，从而有效绕过协议解析器的初步合法性检查。然而，该工具高度依赖专家知识来手动构建 Pit 文件，这在处理复杂或私有协议时存在较高的准入门槛。

Peach*\cite{Luoetal2020} 是在 Peach 基础上改进而成的协议安全测试工具。它继承了 Peach 基于 XML 建模的结构化测试能力，并针对工业控制协议（ICS）的复杂性，引入了覆盖率引导的报文拆解与重组机制。与依赖预定义变异规则的传统 Peach 相比，Peach* 能够实时收集测试过程中的执行路径反馈，将触发新状态的优质报文“碎片化”并重新构建，从而生成更具探测性的测试用例。

Boofuzz\cite{boofuzz} 是目前开源社区中主流的协议测试框架之一。它采用 Python 脚本定义协议结构，提供了强大的会话管理、故障恢复和监控接口。Boofuzz 通过模拟请求与响应的链接关系来探索协议的交互逻辑。相比于 Peach，Boofuzz 更加轻量且易于扩展，但在处理需要覆盖率反馈引导的深度路径探索时，其效率受到盲目变异策略的限制。

AFLNet\cite{pham2020aflnet} 是协议灰盒模糊测试领域的里程碑。它将覆盖率引导的思想引入协议测试，通过监控服务器返回的响应码（如 FTP 的 220 或 SMTP 的 250）来推断协议的有限状态机（FSM）。AFLnet 能够根据状态反馈动态调整种子选择策略，优先探索未触达的协议状态。这种状态感知的机制使得它在针对有状态协议的测试中表现出远超传统工具的路径覆盖能力。

SGFuzz\cite{ba2022sgfuzz} 针对有状态协议中复杂的请求依赖关系提出了状态图引导的变异策略。它通过自动化地构建和利用协议状态图，识别不同消息序列之间的语义依赖性。相比于 AFLnet 随机性较高的序列变异，SGFuzz 能够生成更符合逻辑的消息链，有效减少了因违反协议规约而产生的无效测试用例，从而在复杂的应用层协议测试中实现了更深的漏洞探测。

NSFuzz\cite{qin2023nsfuzz} 关注于网络环境下的反馈同步与测试稳定性问题。由于网络延迟、丢包或超时，模糊测试工具往往难以准确获取目标的真实运行状态。NSFuzz 引入了细粒度的网络状态感知机制，通过优化同步机制和超时处理，降低了网络不确定性带来的反馈干扰。此外，它还强化了对报文内部结构的理解，能够针对网络协议的多样化字段进行更精准的变异，提升了在真实网络拓扑下的测试鲁棒性。


\subsection{大语言模型辅助的协议模糊测试}

越来越多的研究开始探索利用大语言模型来增强协议模糊测试。一些工作尝试使用大语言模型来丰富用于模糊测试的初始种子库，或是从协议规范文档中自动提取协议字段和状态机信息以构建更为准确的协议模型。

ChatAFL \cite{meng2024large} 基于 AFLNet 构建，通过与大语言模型交互提取协议的机器可读语法，从而实现结构感知的变异操作。同时，该方法利用模型生成多样化的消息以扩展初始种子语料库。在模糊测试陷入覆盖率瓶颈时，ChatAFL 进一步通过预测交互序列中的下一条消息，引导测试向新的状态空间探索。

BSFuzzer\cite{yang2026bsfuzzer} 是一个面向 BLE 的规范指导语义模糊测试工具。该工具利用大语言模型从 Bluetooth 规范中提取协议状态机以及报文级语义，并直接合成具有上下文感知能力的测试序列。具体而言，它生成具有语义意义的字段变异以及状态迁移违规行为，将这些操作组合成多步交互序列，并进一步利用基于大语言模型的推理来验证设备响应是否符合预期的协议行为。

MSFuzz \cite{cheng2024msfuzz} 则通过大语言模型从协议实现的源代码中挖掘深入的消息语法知识，而非仅仅依赖于规范文档。其首先利用预处理技术定位关键的消息解析代码，随后通过大语言模型构建精细的消息语法树，以捕获消息类型、字段约束及其合法值。在模糊测试过程中，MSFuzz 利用这些语法知识指导初始种子语料库的扩展，并实现语法感知的变异策略，从而生成更高质量、能通过语法检查的测试用例。

\section{主要研究内容}
纵观协议模糊测试的发展历程，尽管研究重点已从初期的随机变异演进至如今的状态感知与语义启发，但在理解协议字段的深度语义，以及针对这些语义信息进行更有针对性的探索方面，仍存在提升空间。此外，目前学术界的研究重心大多倾向于对以 AFLNet 为代表的基于变异的模糊测试框架进行增强，而对于 Peach 模糊测试框架，相关研究相对较少，且大多集中在为其添加覆盖率反馈上\cite{Luoetal2020,Lian2017}。利用大语言模型的知识抽取与推理能力来解决Peach框架的现有问题，仍是一个尚未充分探索的研究方向。

针对传统 Peach 模糊测试框架协议建模成本高、变异算子缺乏语义感知能力及难以覆盖深层逻辑路径等问题，本文研究并实现了一种大语言模型智能体驱动的协议语义感知变异方法。该方法利用大语言模型的知识抽取与推理能力生成代码，集成到 Peach 的变异流程中，旨在提升测试用例的质量与漏洞发现效率。本文的主要研究内容包括：

\begin{enumerate}[label=\arabic*)]
    \item 设计并实现了一个自动化智能体（Agent）流程，指导大语言模型从协议规范文档中理解语义信息，生成包括数据模型、变异器和修复器在内的协议测试组件；
    \item 设计并实现了与测试组件适配的两阶段变异策略；
    \item 将上述方法集成到 Peach 模糊测试框架中，形成 LLMPeach，并在多个现实世界协议实现上进行了实验评估与分析。
\end{enumerate}

\section{论文组织架构}
第一章\quad 绪论

绪论部分介绍了研究的背景与意义，分析了当前协议模糊测试领域的核心瓶颈与挑战，并概述了国内外在协议模糊测试及大语言模型辅助的协议模糊测试方面的研究现状。最后，明确了本文的主要研究内容和论文的组织架构。

第二章\quad 相关知识

相关知识部分详细介绍了模糊测试的基本概念、Peach 模糊测试框架的核心架构与工作流程，以及大语言模型智能体的原理与应用，为后续章节的研究内容奠定理论基础。

第三章\quad 协议语义感知变异增强方法

协议语义感知变异增强方法部份详细阐述了本文提出的方法的设计，包括大模型智能体驱动的测试组件生成流水线和两阶段变异策略，并介绍了LLMPeach框架的整体架构与实现细节。

第四章\quad 实验评估与分析

实验评估与分析部分通过在多个现实世界协议实现上的实验，评估了所提出的协议语义感知变异增强方法在代码分支覆盖、独特分支等方面的表现，并与基线 Peach 框架进行了对比分析。同时，通过消融研究验证了各个组件对整体性能提升的贡献。

第五章\quad 总结与展望

总结与展望部分回顾了本文的主要研究工作和贡献，分析了研究中存在的不足，并提出了未来可能的研究方向和改进措施。